\chapter{Backend API}

\section{Presentation}

Le backend est une API REST construite avec Express 5 et TypeScript. Son role principal
est de servir de proxy securise vers les APIs NASA, en cachant la cle API cote serveur.

\section{Structure du projet}

\begin{lstlisting}[language={},caption={Arborescence du backend}]
backend/
  src/
    index.ts           # Point d'entree
    app.ts             # Configuration Express
    config/
      index.ts         # Variables d'environnement
      database.ts      # Pool PostgreSQL
      dynamodb.ts      # Client DynamoDB
    routes/
      health.ts        # GET /api/health
      apod.ts          # GET /api/apod, /range, /random
      events.ts        # GET /api/events
      favorites.ts     # CRUD /api/favorites
    services/
      nasaApi.service.ts  # Appels HTTP vers NASA
      cache.service.ts    # Cache DynamoDB
    middleware/
      errorHandler.ts
\end{lstlisting}

\section{Routes API}

\begin{table}[H]
    \centering
    \begin{tabularx}{\textwidth}{llXl}
        \toprule
        \textbf{Methode} & \textbf{Route} & \textbf{Description} & \textbf{Cache} \\
        \midrule
        GET & /api/health & Verification de sante & Non \\
        GET & /api/apod & Image du jour & DynamoDB \\
        GET & /api/apod/range & Images par plage de dates & DynamoDB \\
        GET & /api/apod/random & Images aleatoires & Non \\
        GET & /api/events & Evenements naturels & DynamoDB \\
        GET & /api/favorites & Liste des favoris & Non \\
        POST & /api/favorites & Ajouter un favori & Non \\
        DELETE & /api/favorites/:date & Supprimer un favori & Non \\
        \bottomrule
    \end{tabularx}
    \caption{Endpoints de l'API}
\end{table}

\section{Dockerfile}

Le Dockerfile utilise un build multi-stage pour minimiser la taille de l'image~:

\begin{lstlisting}[language=bash,caption={Dockerfile multi-stage}]
FROM node:20-alpine AS builder
WORKDIR /app
COPY package*.json ./
RUN npm ci
COPY tsconfig.json ./
COPY src/ ./src/
RUN npm run build

FROM node:20-alpine
WORKDIR /app
RUN addgroup -g 1001 appgroup && \
    adduser -u 1001 -G appgroup -s /bin/sh -D appuser
COPY --from=builder /app/dist ./dist
COPY --from=builder /app/package*.json ./
RUN npm ci --omit=dev
USER appuser
EXPOSE 3000
CMD ["node", "dist/index.js"]
\end{lstlisting}

Points cles~:
\begin{itemize}
    \item Utilisateur non-root (\texttt{appuser}) pour la securite
    \item \texttt{npm ci --omit=dev} en production pour reduire la taille
    \item Base \texttt{node:20-alpine} pour une image legere (\~{}180 Mo)
\end{itemize}

\section{Docker Compose --- Developpement local}

\begin{lstlisting}[language=bash,caption={Services docker-compose}]
services:
  backend:     # Express API (port 3000)
  postgres:    # PostgreSQL 16 (port 5432)
  dynamodb-local:  # DynamoDB Local (port 8000)
\end{lstlisting}

Le fichier \texttt{docker-compose.yml} a la racine du projet permet de lancer
l'ensemble de la stack localement avec \texttt{docker compose up}.
